% ============================================================================
\chapter{Programmation fonctionnelle~: lambdas, Streams, collecteurs}
\label{ch-java-fonctionnel}
% ============================================================================

\begin{objectif}
\begin{itemize}
  \item \'Ecrire des lambdas lisibles, comprendre la capture et les interfaces fonctionnelles.
  \item Ma\^itriser le pipeline Streams~: sources, op\'erations interm\'ediaires, terminales.
  \item Utiliser \texttt{Collectors} sans recopier un exemple au hasard~: groupage, partition,
        agr\'egats, maps imbriqu\'ees, collecteurs personnalis\'es.
  \item Savoir quand le flux parall\`ele aide vraiment, et quand il nuit.
  \item \'Ecrire du code fonctionnel qui reste d\'ebogable.
\end{itemize}
\end{objectif}

\section{Lambda~: la syntaxe compl\`ete, puis la r\`egle}

\begin{codejava}
// Les cinq formes, de la plus verbeuse a la plus concise (toutes equivalentes)
Comparator<Integer> c1 = (Integer a, Integer b) -> { return Integer.compare(a, b); };
Comparator<Integer> c2 = (a, b) -> Integer.compare(a, b);
Comparator<Integer> c3 = Comparator.comparingInt(Integer::intValue);
Comparator<Integer> c4 = Integer::compare;
Comparator<Integer> c5 = (a, b) -> a - b;             // correct pour des petits entiers, pas generique !

// Le type cible est infer\'e par le contexte : "var" ne fonctionne pas avec un lambda nu
var r = (Runnable) () -> System.out.println("ok");    // le cast donne le type cible
// Runnable r2 = () -> { };                             // correct
// var r3 = () -> { };                                  // ERREUR de compilation

// Une expression lambda peut etre une expression ou un bloc ; un bloc doit avoir un return.
Runnable vide = () -> { };
java.util.function.Supplier<String> consigne = () -> { var s = "a"; return s + "b"; };

// Les interfaces fonctionnelles du JDK (java.util.function) -- a memoriser par famille :
//   Predicate<T>  : boolean test(T)                 | Function<T,R> : R apply(T)
//   Consumer<T>   : void accept(T)                  | Supplier<T>   : T get()
//   UnaryOperator<T> / BinaryOperator<T>            | BiFunction<T,U,R>
//   IntFunction<R>, ToIntFunction<T>, IntToDoubleFunction, ... (variantes primitives)
//   Runnable, Callable<V> (peut lancer une exception), Comparator<T>, Predicate combinables
Predicate<Integer> pair = i -> i % 2 == 0;
Predicate<Integer> positif = i -> i > 0;
var combo = pair.and(positiv).or(i -> i == 0);
Function<Integer, Integer> fois2 = i -> i * 2;
Function<Integer, Integer> chaine = fois2.andThen(String::valueOf).andThen(s -> "[" + s + "]");
// composition : andThen (apres) vs compose (avant) -- la source d'erreur numero 1

// Capture : uniquement des variables "effectivement finales"
int facteur = 3;
Function<Integer, Integer> mult = x -> x * facteur;      // OK (facteur non modifie ensuite)
// facteur = 4;                                           // -> erreur de compilation : pas de reassignation
final StringBuilder acc = new StringBuilder();            // contournement propre : un objet mutable
Consumer<String> collecter = s -> acc.append(s).append('\n');
// Ce n'est pas une limitation technique mais un choix : capturer une variable locale mutable
// dans un lambda partage entre threads serait une bombe.
\end{codejava}

\begin{note}
En Java, un lambda est une instance d'interface fonctionnelle (pas une fonction de premi\`ere
classe flottante). Le bytecode utilise \verb|invokedynamic| + \verb|LambdaMetafactory|~:
l'objet est cr\'e\'e une fois, et l'appel est aussi rapide qu'une m\'ethode virtuelle. Concr\`etement~:
n'aie pas peur d'un lambda dans une boucle chaude~; aie peur de l'\emph{allocation} d'un
\verb|ArrayList| dans la m\^eme boucle.
\end{note}

\section{Streams~: le pipeline}

Un \verb|Stream| n'est \textbf{pas} une collection~: c'est un plan de calcul sur une source,
\emph{consommable une fois}, paresseux jusqu'\`a l'op\'eration terminale.

\begin{codejava}
Stream<Article> src = articles.stream();

// Intermediaires (paresseux, renvoie un Stream) :
src.filter(a -> a.estPublie() && a.tags().contains("java"))
   .map(Article::titre)
   .map(String::trim)
   .distinct()
   .sorted(Comparator.reverseOrder())
   .peek(t -> log.info("traitement de {}", t))        // peek = debug seulement
   .limit(20)
   .skip(5)
   .flatMap(t -> Stream.of(t.split(" ")))              // un element -> plusieurs
   .mapMulti(String.class, (t, out) -> { for (var w : t.split(" ")) { out.accept(w); } });

// Terminales (\'evaluent tout, une seule fois) :
List<String> l = stream.toList();                       // Java 16+ (immuable)
var copieModifiable = stream.collect(Collectors.toCollection(ArrayList::new));
Set<String> set = stream.collect(Collectors.toSet());
Map<Long, String> map = stream.collect(Collectors.toMap(Article::id, Article::titre));
Optional<Article> premier = stream.findFirst();          // rencontre le 1er (ordre = source)
Article un = stream.findAny().orElseThrow();              // en parallele : findAny est moins restrictif
long n = stream.count();
boolean aucun = stream.noneMatch(pred);                  // anyMatch/allMatch/noneMatch = courts-circuits
OptionalInt min = stream.mapToInt(Article::vues).min();
int somme = stream.map(Article::vues).reduce(0, Integer::sum);
String joint = stream.map(Article::titre).collect(Collectors.joining(", ", "[", "]"));
stream.forEach(System.out::println);                     // pas forEachOrdered (garde l'ordre en parallele)
var tableau = stream.toArray(String[]::new);
var optionnel = Optional.ofNullable(stream.findFirst().orElse(null));
\end{codejava}

\subsection{Les collecteurs, l\`a o\`u le talent se voit}

\begin{codejava}
import java.util.stream.Collectors;
import static java.util.stream.Collectors.*;

// Agr\'egations
long nbPublies = articles.stream().filter(Article::estPublie).count();
DoubleSummaryStatistics stats = articles.stream().mapToDouble(Article::vues).summaryStatistics();
stats.getCount(); stats.getAverage(); stats.getMin(); stats.getMax(); stats.getSum();
IntSummaryStatistics notes = eleves.stream().mapToInt(Eleve::note).summaryStatistics();

// Groupages
Map<Statut, List<Article>> parStatut = articles.stream().collect(groupingBy(Article::statut));
Map<Statut, Long> compteParStatut = articles.stream().collect(groupingBy(Article::statut, counting()));
Map<String, List<String>> titresParAnnee = articles.stream().collect(
        groupingBy(a -> String.valueOf(a.publieLe().getYear()),
                   collectingAndThen(mapping(Article::titre, toList()), l -> new ArrayList<>(l))));
Map<Statut, Map<String, List<Article>>> doubleCle = articles.stream()
        .collect(groupingBy(Article::statut, groupingBy(a -> a.tags().isEmpty() ? "sans-tag" : a.tags().get(0))));
Map<Boolean, List<Article>> partition = articles.stream().collect(partitioningBy(Article::estPublie));

// Maps et transformations
Map<Long, Article> parId = articles.stream().collect(toMap(Article::id, a -> a));
Map<Long, Article> parIdOrdonne = articles.stream().collect(toMap(Article::id, a -> a, (x, y) -> x, TreeMap::new));
Map<String, Integer> totauxParTag = articles.stream()
        .flatMap(a -> a.tags().stream())
        .collect(groupingBy(t -> t, summingInt(t -> 1)));
// (plus simple : groupingBy + counting)

// Collections cibles
var parSet = articles.stream().collect(toCollection(LinkedHashSet::new));
String csv = articles.stream().map(Article::titre).collect(joining(",", "\"", "\""));
Map<String, Long> top5 = stream.collect(groupingBy(f, counting())).entrySet().stream()
        .sorted(Map.Entry.<String, Long>comparingByValue().reversed()).limit(5)
        .collect(toMap(Map.Entry::getKey, Map.Entry::getValue, (a, b) -> a, LinkedHashMap::new));

// Collecteur sur mesure (reduction mutable propre)
record Compacteur(int total, int max, List<String> noms) {
    static Compacteur creer() { return new Compacteur(0, 0, new ArrayList<>()); }
    void ajouter(String nom, int valeur) { noms.add(nom); total += valeur; max = Math.max(max, valeur); }
    Compacteur fusionner(Compacteur o) {   // utilise en parallele
        noms.addAll(o.noms); return new Compacteur(total + o.total, Math.max(max, o.max), noms);
    }
}
var resume = flux.collect(Collector.of(
        Compacteur::creer,
        (c, a) -> c.ajouter(a.titre(), a.vues()),
        Compacteur::fusionner,
        Collector.Characteristics.CONCURRENT));

// Un stream ne se relit pas : si tu as besoin du resultat deux fois, materialise-le.
List<Article> materialise = articles.stream().filter(c::pertinent).toList();
\end{codejava}

\begin{tabz}{lXl}
  \headrow \thd{Collecteur} \& \thd{\`A la place de...} \& \thd{Note} \\
  \midrule
  \verb|toList()| (Java~16+) \& \verb|collect(toList())| \& immuable~; \verb|collect(Collectors.toList())| reste mutable. \\
  \verb|toUnmodifiableList()| \& -- \& Java~10+, et \'echoue si un \'el\'ement est \verb|null|. \\
  \verb|toMap(k, v)| \& une boucle + put | \textbf{l\`eve} sur cl\'e dupliqu\'ee~: passe le mergeur \verb|(a,b)->a|. \\
  \verb|toCollection(HashSet::new)| \& -- \& type de collection explicite. \\
  \verb|groupingBy(classifier)| \& \verb|Map<K, List<V>>| rempli \`a la main \& accepte un downstream. \\
  \verb|groupingBy(classifier, downstream)| \& -- \& \verb|counting|, \verb|summingInt|, \verb|reducing|, \verb|minByBy|... \\
  \verb|partitioningBy(pred)| \& deux listes \& 2 cl\'es seulement, plus rapide qu'un groupage. \\
  \verb|joining(delim, prefix, suffix)| \& \verb|StringBuilder| \& attention : lent sur de gros volumes. \\
  \verb|mapping(f, d)|, \verb|filtering(p, d)|, \verb|collectingAndThen(...)| \& -- \& la combinatoire qui \'evite un stream interm\'ediaire. \\
  \verb|summarizingInt/Long/Double| \& trois parcours | un seul. \\
  \verb|reducing(id, f, op)| \& un accumulateur | associativit\'e obligatoire pour le parall\`ele. \\
\end{tabz}

\section{Flux primitifs~: la performance sans emballage}

\begin{codejava}
// Stream<Integer> cree des objets Integer pour chaque element. Les flux primitifs non.
int total = articles.stream().mapToInt(Article::vues).sum();
long totalLong = articles.stream().mapToLong(Article::vues).asLongStream().sum();
double moyenne = articles.stream().mapToDouble(a -> a.vues()).average().orElse(0);
DoubleStream.of(1.1, 2.2).map(Math::sqrt).filter(d -> d > 1).summaryStatistics();
IntStream.range(0, n).map(i -> i * i).filter(i -> i % 2 == 0).boxed().toList();
IntStream.rangeClosed(1, 100).sum();                      // 1..100 (range = exclut la fin)
IntStream.iterate(1, v -> v <= 1000, v -> v * 2).toArray();
IntStream.generate(() -> (int) (Math.random() * 6) + 1).limit(10).toArray();
Stream.generate(() -> "constant").limit(3);
Stream.iterate(new int[]{0, 1}, t -> new int[]{t[1], t[0] + t[1]}).limit(10).map(t -> t[0]);  // Fibonacci
// Boite / deboite : .boxed() (primitif -> objet) ; .mapToInt() (objet -> primitif)
List<Integer> entiers = IntStream.of(1, 2, 3).boxed().toList();
OptionalInt vide = IntStream.empty().min();
int avecDefaut = IntStream.of().max().orElse(-1);        // les Optional primitifs ont orElse/orElseGet
\end{codejava}

\section{Parall\'elisme de flux~: le mode d'emploi honn\^ete}

\begin{codejava}
int somme = fluxEnMemoire.parallel().mapToInt(String::length).sum();

// Le parallel() utilise le FORK/JOIN COMMUN (threads = nb de coeurs - 1). Donc :
//  - interdit d'y faire du bloquant (IO, sommeil, appel HTTP) : tu affames les autres ;
//  - inutile sur une liste de 5 000 elements : le cout de partition > le gain ;
//  - rentable sur des donnees en memoire, sans \'etat partag\'e, avec une operation courte et pure ;
//  - a mesurer : parfois 8x, parfois 0.9x (contention, allocation, ordre impos\'e).
List<String> ok = stream.parallel().map(String::trim).filter(s -> !s.isEmpty()).toList();
// Operations "sans ordre" (findAny, noneMatch, unordered) sont les amies du parallele :
Optional<String> nImporteLequel = stream.parallel().unordered().filter(c::estBon).findAny();
// Limiter le pool pour un calcul isol\'e, sans polluer le commun :
try (var pool = new ForkJoinPool(4)) {                  // Java 19+: AutoCloseable
    resultat = pool.submit(() -> stream.parallel().map(...).reduce(...)).get();
}
// Regle d'or : si un calcul est trop gros pour un thread et trop petit pour un cluster,
// un stream parallele est la bonne idee. Sinon, les virtual threads (chapitre suivant) sont
// presque toujours meilleurs.
\end{codejava}

\subsection{\'Ecrire du code fonctionnel qui reste lisible}

\begin{enumerate}
  \item Une op\'eration par ligne dans les pipelines de plus de trois \'etapes~; retour \`a la ligne
        apr\`es chaque \verb|.|, indent\'es de quatre espaces.
  \item Un lamba de plus de trois lignes~: nomme-le (m\'ethode priv\'ee)~: \verb|a -> estPublieRecent(a)|.
  \item \verb|peek| n'est pas un outil de production~: il peut \^etre ignor\'e par le moteur
        (court-circuit). Utilise un compteur ou un vrai point de log.
  \item Pas d'effet de bord dans \verb|map|/ \verb|filter|~: le compilateur ne te sauvera pas,
        mais le parall\`ele et le r\'eordonnancement, si.
  \item Un stream consomm\'e est mort~: si tu it\`eres deux fois, mat\'erialise une \verb|List|.
  \item Ferme les flux issus de \verb|Files.lines| (try-with-resources)~: ils d\'etiennent un handle.
  \item La boucle \verb|for| n'est pas un aveu d'\'echec~: parcours avec indices, mutation d'un
        tableau, algorithme avec break conditionnel~? La boucle gagne. Les streams excellent sur
        \og ~transformer une collection en une autre~ \fg{}.
\end{enumerate}

\begin{codejava}
// Refactorisation type : de "Java 5" a "Java 21", en gardant la clarte
// AVANT
Map<String, List<Article>> res = new HashMap<>();
for (Article a : articles) {
    if (a.estPublie()) {
        for (String tag : a.tags()) {
            res.computeIfAbsent(tag, k -> new ArrayList<>()).add(a);
        }
    }
}
// APRES
Map<String, List<Article>> res = articles.stream()
        .filter(Article::estPublie)
        .flatMap(a -> a.tags().stream().map(t -> Map.entry(t, a)))
        .collect(groupingBy(Map.Entry::getKey, mapping(Map.Entry::getValue, toList())));
// et un jour, tu realise que computeIfAbsent etait plus rapide a lire : tu reviens en arriere.
// La bonne question n'est pas "fonctionnel ou boucle" mais "le lecteur comprend-il en 5 secondes".
\end{codejava}

\begin{exercice}
\begin{enumerate}
  \item \`A partir d'une liste de 100 000 enregistrements \verb|Vente(id, produit, region, montant, date)|,
        produis~: le chiffre par r\'egion~; le top 10 produits~: la croissance mensuelle~; le panier
        moyen~; les r\'egions sans vente~; et un export CSV tri\'e~: le tout en streams~: puis en boucles~;
        compare la lisibilit\'e, pas seulement le temps.
  \item \'Ecris un collecteur \verb|versStatistiques()| qui renvoie un \verb|record Stats(n, somme,
        moyenne, min, max, ecartType)| en un seul passage~: puis valide-le avec des propri\'et\'es
        (valeurs al\'eatoires, comparaison avec trois parcours s\'epar\'es).
  \item Construis un mini-\verb|Collectors| : \verb|toMapMulti|, \verb|toGroupedCount|, \verb|toSortedSet|.
  \item Utilise \verb|mapMulti| pour aplatisir un arbre (\verb|record Noeud(String nom, List<Noeud> enfants)|)
        en un flux de chemins \og ~a/b/c~ \fg{}~: puis compare \`a un \verb|flatMap| r\'ecursif.
  \item Calcule les premiers 10 000 nombres premiers~: (a) filtre na\"if en stream~; (b) crible
        d'\'Eratosth\`ene avec \verb|int[]|~; (c) version parall\`ele. Donne les trois temps et explique
        l'\'ecart. (Spoiler~: le stream sera 50 fois plus lent, et c'est formateur.)
  \item Ajoute une m\'ethode par d\'efaut \`a une de tes interfaces fonctionnelles~: \verb| puis |,
        \verb|enSecurite| (capture les exceptions~: \verb|Function<T,R>| -> \verb|Resultat<R>|)~;
        utilise-la pour transformer un parsing fragile en composition propre.
\end{enumerate}
\end{exercice}

\begin{ressources}
\begin{itemize}
  \item Tutorial officiel Streams --- \url{https://dev.java/learn/streams/}
  \item Javadoc \verb|Collectors| (le tableau des collecteurs vaut le d\'etour) --- \url{https://docs.oracle.com/en/java/javase/21/docs/api/java.base/java/util/stream/Collectors.html}
  \item \emph{Java Streams} (P. Sotuyodhy, en fran\c{c}ais) --- \url{https://blog.sotuyodhy.com/java-streams/}
  \item Article de Brian Goetz sur les limites du parall\'elisme --- \url{https://www.oracle.com/technical-resources/articles/java/architect-streams.html}
  \item Vavr (fonctionnel avanc\'e en Java) --- \url{https://www.vavr.io/}
\end{itemize}
\end{ressources}

\index{lambda}\index{Streams}\index{Collectors}\index{parallel stream}
